\paragraph{边界点处的赋值。}\label{sub:evaluation-boundary-points}
设给定了一个带\trLagrange\trLabel{}的\trPointedBoundary\trRiemann\trSurface，外加一个有限的%
\trMarkedPoint{}集合 $\Pi\subset\partial S$，分成正的与负的两组 $\Pi^-$、$\Pi^+$。于是便可
对 TQFT 映射\eqref{eq:tqft-operation} 引入一个推广：
\begin{equation}\label{eq:evaluation-tqft-map}
\begin{aligned}
\Phi_{S,\Pi}\colon
&\bigotimes_{\zeta^+\in\Sigma^+}HF^{\mathrm{pr}}(L_{\zeta+,0},L_{\zeta+,1})
\otimes
\bigotimes_{z^+\in\Pi^+}H^*(L_{C_{z+}};\KK)\\
&\longrightarrow
\bigotimes_{\zeta^-\in\Sigma^-}HF^{\mathrm{pr}}(L_{\zeta-,0},L_{\zeta-,1})
\otimes
\bigotimes_{z^-\in\Pi^-}H^*(L_{C_{z-}};\KK).
\end{aligned}
\end{equation}
这里 $C_z$ 是 $z$ 所属的那个边界\trConnectedComponent。除了通常的把无穷远点%
\trGluing{}起来之外，现在还可以作\emph{\trBoundaryConnectedSum}（boundary
connected sums），它把两个\trSurface{}$S_1$、$S_2$ 在点 $z_1^+\in\Pi_1^+$、
$z_2^-\in\Pi_2^-$ 处熔接起来，那里各自的\trLagrange\trBoundaryCondition{}相互重合。
对与之相关联的算子\eqref{eq:evaluation-tqft-map}，这对应于把\trComposition{}作
用到相关的因子 $H^*(L_{C_z};\KK)$ 上。
最简单的情形是\trCompact{}的 $S$。假设\trRegularity{}成立，则%
\eqref{eq:inhomogeneous-equation} 之解的\trModuliSpace{}$\scr{M}_S$ 将是一个%
\trSmooth\trClosed\trManifold。于是\trEvaluationMap{}$\mathrm{ev}_z(u)=u(z)$（$z\in\Pi$）的
乘积就给出一个落在 $\prod_z L_{C_z}$ 中的\trHomology{}类。人们用%
\trKunnethFormula{}与\trPoincareDuality{}（对 $\Pi^-$ 因子）把它读成一个%
\trMultilinear{}映射\eqref{eq:evaluation-tqft-map}。然而，若把\trInhomogeneous{}项
取为零，则一个直截了当的\trEnergy{}论证表明 $\scr{M}_S$ 只含常值映射，这意味
着所得的\trInvariant{}其实没有多大意思。例如，取 $S=D$ 配以\trLagrange\trLabel{}$L$，
就得到
\begin{equation}\label{eq:disc-evaluation-invariants}
\begin{cases}
\text{单位元 }1\in H^*(L;\KK), & |\Pi^-|=1,\ |\Pi^+|=0,\\
\text{积分映射 }H^*(L;\KK)\to\KK, & |\Pi^-|=0,\ |\Pi^+|=1,\\
\text{恒等映射 }H^*(L;\KK)\to H^*(L;\KK), & |\Pi^-|=|\Pi^+|=1,\\
\text{\trCupProduct{}（cup product），} & |\Pi^-|=1,\ |\Pi^+|=2.
\end{cases}
\end{equation}
为对更一般的情形给一点简短的提示，假设 $\Sigma=\{\zeta\}$ 由单个出向端组成，
而 $\Pi=\{z\}$ 由单个负\trMarkedPoint{}组成。其想法是：若取一个\trFloer\trCocycle
$\sum\limits_k m_k y_k\in CF^{\mathrm{pr}}(L_{\zeta,0},L_{\zeta,1})$，则相应的、由%
\trEvaluationMap{}$\mathrm{ev}_z\colon\scr{M}(y_k)\to L_{C_z}$ 组成的形式\trLinear{}组合是
$L_z$ 上普通\trTopology{}意义下的一个\trCycle，其类（或者更确切地说，它的%
\trPoincareDuality）将是\trFloerCohomology{}类 $[\sum\limits_k m_k y_k]$ 在%
\eqref{eq:evaluation-tqft-map} 之下的像。一个与之\trEquivalence{}、但在技术上或许
更整洁的途径
（它可追溯到 [120]）是为 $H^*(L_{C_z};\KK)$ 采用\trMorseHomology{}模型。这意味
着，处于\eqref{eq:evaluation-tqft-map} 底层的上\trChainMap，其\trCoefficients{}由计数交
$\scr{M}(y_k)\cap\mathrm{ev}_z^{-1}(W^u(x))$ 中的点来确定，其中 $W^u(x)$ 是
$-\nabla h_z$ 的\trUnstableManifold，而 $h_z$ 是取定在 $L_{C_z}$ 上的一个%
\trMorseFunction。例如，取 $S=\overline{H}$ 配以\trLagrange\trLabel{}$L$，便定义了
PSS 映射\eqref{eq:pss-map}。\cref{sub:donaldson-fukaya-category} 中陈述的一切
性质现在都可由\eqref{eq:disc-evaluation-invariants} 与\trCompositionLaw{}推出。
例如，其逆映射恰好由带一个正\trMarkedPoint{}的 $S=H$ 给出。在某种意义上，
这其实是唯一的应用：人们可以借助\trGluing{}与 PSS \trIsomorphism，把%
\trMarkedPoint{}$z$ 换成满足
$L_{\zeta,0}=L_{\zeta,1}$ 的无穷远点 $\zeta$，这就把\eqref{eq:evaluation-tqft-map}
化归为更基本的版本\eqref{eq:tqft-operation}。
